\documentclass[12pt]{beamer}
\usetheme{Singapore}

\usepackage[utf8]{inputenc}
\usepackage{tikz}
%\usetikzlibrary{patterns}
\usetikzlibrary{shapes.geometric}
\usetikzlibrary{calc}
%\usetikzlibrary{automata,positioning}
\usepackage{siunitx}
\usepackage[e]{esvect}

\usepackage{array}
\usepackage{chemfig}
\usepackage{amsmath}
\usepackage{amssymb}
%\usepackage[cal=boondoxo]{mathalfa} 
%\usepackage{enumitem}
%\usepackage{lastpage}
\usepackage{tcolorbox}
%\usepackage{siunitx}

\renewcommand{\arraystretch}{1.5}
\setbeamersize{text margin left=3mm,text margin right=3mm}


\begin{document}
\title{Physique Chimie en seconde}
\author{ROSALIE}
\date{2024 - 2025}

\begin{frame}
\titlepage
\end{frame}

\begin{frame}
    \frametitle{Matériel à apporter}
    \begin{itemize}
        \item Classeur de physique-chimie
        \item Calculatrice (lycée ou collège)
        \item Blouse en coton pour les TP de chimie
        \item Trousse complète, critérium, crayons de couleur
    \end{itemize}
\end{frame}

\begin{frame}
    \frametitle{Manuel choisi : Le livre scolaire}
    La version numérique, gratuite, sans inscription :\\
    https://www.calameo.com/read/0005967298a6f595a36e9

\end{frame}

\begin{frame}
    \begin{center}
        \huge À quoi servent la Physique et la Chimie ?
    \end{center}
\end{frame}

\begin{frame}
    \begin{center}
        \huge À expliquer et prévoir le comportement de ce qui n'est pas vivant.
    \end{center}
\end{frame}

\begin{frame}
    \begin{center}
        \huge Comment parvient-on à le faire ?
    \end{center}
\end{frame}

\begin{frame}
    \begin{center}
        \huge Par des expériences !
    \end{center}
    Les résultats des expériences permettent par exemple :
    \begin{itemize}
        \item d'imaginer de nouveaux modèles théoriques
        \item de valider ou d'invalider un modèle théorique
    \end{itemize}
\end{frame}

\begin{frame}
    \frametitle{Programme de l'année}
    \begin{itemize}
        \item Thème 1 : Constitution et transformation de la matière
        \item Thème 2 : Mouvements et interactions
        \item Thème 3 : Ondes et signaux
    \end{itemize}
\end{frame}

\begin{frame}
    \frametitle{Les évaluations}
    Tout au long de l'année :
    \begin{itemize}
        \item Activités ou TP pourront être évaluées (coefficient 1)
        \item Devoirs bilans (durée maximale 1h) seront programmées.\\
        \begin{itemize}
            \item contenu : 2 chapitres en général
            \item fréquence par trimestre 2 ou 3
            \item coefficient  3
        \end{itemize}
    \end{itemize}
\end{frame}
\end{document}